\labsection{1}{Programming foundations and structured records}{sec:lab01}

\begin{sourcealignment}{Official Lab Manual --- Lab I}
\textbf{Objective.} Revise programming prerequisites needed later in the course. The manual explicitly names structures, unions, pointers, and 1D arrays.

\textbf{Observed outcomes.} Understand data structures and their types; use data structures for unseen problems; select an appropriate structure for a real-world problem.

\textbf{Problem directions.} Basic C/C++ programs using loops and structures; programs on unions and 1D arrays; a small database-style application using structure, pointer, and array.
\end{sourcealignment}

\begin{labproject}{Build a small record-based application}
\textbf{Coverage.} Chapter 1.

The lecture task asks for at least five student records containing ID, name, age, and course, followed by display, search-by-ID, and update operations. This also satisfies the lab manual's database-style direction.

\begin{lstlisting}
#include <iostream>
#include <string>
using namespace std;

struct Student {
    int id;
    string name;
    int age;
    string course;
};

void display(const Student students[], int n) {
    for (int i = 0; i < n; ++i) {
        cout << students[i].id << " | "
             << students[i].name << " | "
             << students[i].age << " | "
             << students[i].course << '\n';
    }
}

int findById(const Student students[], int n, int id) {
    for (int i = 0; i < n; ++i)
        if (students[i].id == id) return i;
    return -1;
}
\end{lstlisting}

\begin{tryit}
Add \texttt{updateStudent()} and use a pointer to the selected record once the matching index is known. Then add one small 1D-array statistic such as the average student age.
\end{tryit}

\begin{labmode}
The manual names \emph{union} as prerequisite review but does not provide detailed union theory in the supplied text. Treat union exercises as revision work directed by your instructor rather than as a new data-structures topic.
\end{labmode}
\end{labproject}

\begin{verification}
\begin{itemize}
  \item At least five records are stored and displayed correctly.
  \item A valid ID is found; an invalid ID returns a clear ``not found'' result.
  \item Updating one record changes only the intended record.
  \item Array bounds are respected.
  \item The role of a pointer versus an array index can be explained.
\end{itemize}
\end{verification}

\begin{labdeliverable}
Submit the C++ program, sample output for successful/unsuccessful search, and a short note identifying the data structure used for the collection and the structure used for each record.
\end{labdeliverable}
